\documentclass{article}
\usepackage[utf8]{inputenc}
\usepackage[russian]{babel}
\usepackage{amsmath}
\usepackage{amssymb}
\usepackage{mathtools}
\usepackage[a4paper, left=1.5cm, right=1.5cm, top=1.5cm, bottom=1.5cm]{geometry}\setlength{\parindent}{0pt}

\title{ Билет 7. Определение равномерной сходимости. Критерий Коши. Теоремы об интегрируемости и
непрерывности предела равномерно сходящейся функции.}
\author{}
\date{}

\begin{document}

\maketitle

\subsection*{Определение равномерной сходимости}

Рассмотрим $f(x, y)$, где $y \in Y \subset \mathbb{R}^1$, и при каждом фиксированном $\bar{y}$ определен интеграл $\int_{\alpha(y)}^{\beta(y)} f(x, y) \, dx$.

Теперь рассмотрим общий случай. Пусть у нас есть функция $f(x, y)$, где $(x, y) \in X \times Y$.

Для каждого фиксированного $x \in X$ $\lim_{y \to y_0} f(x, y)$ предел существует и конечен, который мы обозначим как $\varphi(x)$. То есть,
\[
\forall x \in X : \quad \lim_{y \to y_0} f(x, y) = \varphi(x).
\]
Это означает, что функция $f(x, y)$, рассматриваемая как функция переменной $y$ при фиксированном $x$, имеет предел $\varphi(x)$ при $y \to y_0$.

Тогда говорят, что на множестве $X$ функция $\varphi(x)$ является \textit{предельной функцией} для $f(x, y)$ при $y \to y_0$.

Записывают это так:
\[
f(x, y) \xrightarrow[y \to y_0]{} \varphi(x) \quad \text{(поточечно на } X \text{ сходится)}.
\]

\textbf{Введем понятие равномерной сходимости.}

Говорят, что $f(x, y)$ стремится к $\varphi(x)$ \textit{равномерно по} $x \in X$ при $y \to y_0$, если:

\[
\boxed{
\forall \varepsilon > 0 \quad \exists \delta > 0 : \forall y \in Y, \; 0 < |y - y_0| < \delta, \quad \forall x \in X \quad \text{выполняется} \quad |f(x, y) - \varphi(x)| < \varepsilon.
}
\]

Это означает, что $\delta$ зависит только от $\varepsilon$, но \textbf{не зависит} от $x$. Это ключевое отличие от поточечной сходимости.

Запись: \quad $f(x, y) \mathop{\rightrightarrows}\limits^{ x \in X}_{\,y \to y_0\,} \varphi(x)$.


\subsection*{Пример (на поточечную, но не равномерную сходимость)}

Рассмотрим $f(x, y) = xy$, где $x \in \mathbb{R}^1, y \in \mathbb{R}^1$, и $y_0 = 0$.

Тогда $\lim_{y \to 0} f(x, y) = \varphi(x) \equiv 0$.

Покажем, что сходимость \textbf{не является равномерной} на $\mathbb{R}$.

Для этого возьмем $\varepsilon = 1$. Тогда для любого $\delta > 0$ можно выбрать $y = \frac{\delta}{2}$ и $x = \frac{2}{\delta}$. Тогда:
\[
|f(x, y) - \varphi(x)| = \left| \frac{2}{\delta} \cdot \frac{\delta}{2} \right| = 1 = \varepsilon.
\]
Следовательно, $\exists \varepsilon > 0$ такое, что $\forall \delta > 0$ найдутся $x, y$ с $0 < |y - y_0| < \delta$, но $|f(x, y) - \varphi(x)| \geq \varepsilon$. Значит, сходимость не равномерна.

\textbf{Однако}, если ограничить $x$: например, рассмотреть $|x| \leq 1$, то сходимость будет равномерной. Действительно, для $\varepsilon > 0$ выберем $\delta = \varepsilon$. Тогда при $0 < |y| < \delta$ и $|x| \leq 1$ имеем:
\[
|f(x, y) - \varphi(x)| = |xy| \leq |x||y| < 1 \cdot \varepsilon = \varepsilon.
\]

\subsection*{Критерий Коши для равномерной сходимости}

\textbf{Условие Коши}: Для того чтобы \quad $f(x, y) \mathop{\rightrightarrows}\limits^{x \in X}_{\,y \to y_0\,} \varphi(x)$, необходимо и достаточно, чтобы

\[
\boxed{\forall \varepsilon > 0 \quad \exists \delta > 0 : \forall y', y'' \in Y, \; 0 < |y' - y_0| < \delta, \; 0 < |y'' - y_0| < \delta, \quad \forall x \in X,}
\]
\[
\boxed{|f(x, y') - f(x, y'')| < \varepsilon.}
\]

\textbf{Доказательство (необходимость)}:

Предположим, что $f(x, y) \to \varphi(x)$ равномерно по $x \in X$ при $y \to y_0$.  
По определению равномерной сходимости:
\[
\forall \varepsilon > 0\ \exists \delta > 0\ \forall y \in Y,\ 0 < |y - y_0| < \delta,\ \forall x \in X:\quad 
|f(x, y) - \varphi(x)| < \frac{\varepsilon}{2}.
\tag{1}
\]

Возьмём произвольные $y', y'' \in Y$, удовлетворяющие $0 < |y' - y_0| < \delta$ и $0 < |y'' - y_0| < \delta$, и фиксируем произвольное $x \in X$.  
Тогда, применяя неравенство треугольника, получаем:
\[
\begin{aligned}
|f(x, y') - f(x, y'')|
&= \bigl| (f(x, y') - \varphi(x)) + (\varphi(x) - f(x, y'')) \bigr| \\
&\le |f(x, y') - \varphi(x)| + |f(x, y'') - \varphi(x)| \\
&< \frac{\varepsilon}{2} + \frac{\varepsilon}{2} = \varepsilon.
\end{aligned}
\]
Так как $x \in X$ произвольно, условие (К) выполнено.

\textbf{Доказательство (достаточность)}:

Пусть выполнено условие (К). Покажем, что существует функция $\varphi \colon X \to \mathbb{R}$, такая что $f(x, y) \to \varphi(x)$ равномерно по $x \in X$.

\textit{Шаг 1: построение $\varphi(x)$.}  
Зафиксируем произвольное $x \in X$. Условие (К), применённое к этому $x$, означает, что числовая функция $y \mapsto f(x, y)$ удовлетворяет критерию Коши при $y \to y_0$. Следовательно, существует конечный предел:
\[
\varphi(x) := \lim_{y \to y_0} f(x, y).
\]
Этим определена функция $\varphi \colon X \to \mathbb{R}$.

\textit{Шаг 2: доказательство равномерной сходимости.}  
Зафиксируем $\varepsilon > 0$. По условию (К) найдётся $\delta > 0$, такое что для всех $y', y'' \in Y$ с $0 < |y' - y_0| < \delta$, $0 < |y'' - y_0| < \delta$ и всех $x \in X$:
\[
|f(x, y') - f(x, y'')| < \frac{\varepsilon}{2}.
\tag{2}
\]

Теперь зафиксируем произвольное $y \in Y$, удовлетворяющее $0 < |y - y_0| < \delta$, и произвольное $x \in X$. Поскольку $\varphi(x) = \lim_{y'' \to y_0} f(x, y'')$, существует последовательность $(y_n) \subset Y$, $y_n \ne y_0$, $y_n \to y_0$, такая что $f(x, y_n) \to \varphi(x)$. Выберем $n$ достаточно большим, чтобы $0 < |y_n - y_0| < \delta$. Тогда из (2) следует:
\[
|f(x, y) - f(x, y_n)| < \frac{\varepsilon}{2}.
\]
Переходя к пределу при $n \to \infty$, с учётом непрерывности модуля, получаем:
\[
|f(x, y) - \varphi(x)| = \lim_{n \to \infty} |f(x, y) - f(x, y_n)| \le \frac{\varepsilon}{2} < \varepsilon.
\]

Поскольку $x \in X$ и $y$ (с условием $0 < |y - y_0| < \delta$) выбирались произвольно, мы доказали, что:
\[
\forall \varepsilon > 0\ \exists \delta > 0\ \forall y \in Y,\ 0 < |y - y_0| < \delta,\ \forall x \in X:\quad |f(x, y) - \varphi(x)| < \varepsilon,
\]
то есть $f(x, y) \to \varphi(x)$ равномерно по $x \in X$.

\subsection*{Теорема о непрерывности предела равномерно сходящейся функции}

Пусть $f(x, y)$ определена на $[a, b] \times Y$, и $\forall y \in Y$ функция $f(x, y)$ непрерывна по $x$ на $[a, b]$.

Если $\forall y \in Y$: $f(x, y) \mathop{\rightrightarrows}\limits^{x \in [a,b]}_{\,y \to y_0\,} \varphi(x)$,
то предельная функция тоже непрерывна на $[a, b]$.


\textbf{Доказательство}:

Покажем, что $\varphi(x)$ непрерывна в произвольной точке $x_0 \in [a, b]$.

$\forall \varepsilon > 0$ $\exists \delta > 0$ : $\forall x \in [a, b]$, $|x - x_0| < \delta$.
Следовательно.
\[
|f(x, y) - \varphi(x)| < \frac{\varepsilon}{3}.
\]

Также, поскольку $f(x, y)$ непрерывна по $x$ для каждого фиксированного $y$, то для выбранного $y$ (например, зафиксируем $y$ с $0 < |y - y_0| < \delta_1$) $\exists \delta_2 > 0$ такое, что $\forall x \in [a, b]$ с $|x - x_0| < \delta_2$ выполняется:
\[
|f(x, y) - f(x_0, y)| < \frac{\varepsilon}{3}.
\]

Теперь для $x$ с $|x - x_0| < \delta_2$ оценим $|\varphi(x) - \varphi(x_0)|$:
\[
|\varphi(x) - \varphi(x_0)| \leq |\varphi(x) - f(x, y)| + |f(x, y) - f(x_0, y)| + |f(x_0, y) - \varphi(x_0)| < \frac{\varepsilon}{3} + \frac{\varepsilon}{3} + \frac{\varepsilon}{3} = \varepsilon.
\]

Таким образом, $\varphi(x)$ непрерывна в точке $x_0$. Поскольку $x_0$ — произвольная точка из $[a, b]$, то $\varphi(x)$ непрерывна на всем отрезке.

\textbf{Замечание из вашего конспекта:} Если $\forall y \in Y$ функция $f(x, y)$ непрерывна по $x$, то предельная функция $\varphi(x)$ тоже может быть непрерывной, но это гарантировано только при равномерной сходимости.


\subsection*{Теорема об интегрируемости предела равномерно сходящейся функции}

Пусть $f(x, y)$ определена на $(x, y) \in [a, b] \times Y$, и $\forall y \in Y$ существует интеграл
\[
\int_a^b f(x, y) \, dx.
\]

Если $f(x, y) \mathop{\rightrightarrows}\limits^{x \in [a,b]}_{\,y \to y_0\,} \varphi(x)$ , то существует интеграл
\[
\int_a^b \varphi(x) \, dx,
\]

\textbf{Доказательство}:

1. \textbf{Интегрируемость $\varphi(x)$}.

По определению равномерной сходимости: $\forall \varepsilon > 0$ $\exists \delta > 0$ такое, что $\forall y \in Y$ с $0 < |y - y_0| < \delta$ и $\forall x \in [a, b]$ выполняется:
\[
|f(x, y) - \varphi(x)| < \frac{\varepsilon}{4(b-a)}.
\]

Выберем конкретное $y$ с $0 < |y - y_0| < \delta$. Поскольку $f(x, y)$ интегрируема, то для данного $\varepsilon > 0$ $\exists$ разбиение $\tau$ отрезка $[a, b]$ на части $[x_i, x_{i+1}]$ такое, что сумма колебаний $\sum_{i=0}^{n-1} \omega_i \Delta x_i < \frac{\varepsilon}{2}$, где $\omega_i = M_i - m_i$ — колебание $f(x, y)$ на $[x_i, x_{i+1}]$, а $\Delta x_i = x_{i+1} - x_i$.

Рассмотрим колебание $\varphi(x)$ на том же отрезке $[x_i, x_{i+1}]$. Пусть $M_i^\varphi = \sup_{x \in [x_i, x_{i+1}]} \varphi(x)$, $m_i^\varphi = \inf_{x \in [x_i, x_{i+1}]} \varphi(x)$.

Тогда для любого $x \in [x_i, x_{i+1}]$:
\[
f(x, y) - \frac{\varepsilon}{4(b-a)} < \varphi(x) < f(x, y) + \frac{\varepsilon}{4(b-a)}.
\]

Отсюда следует, что
\[
M_i^\varphi \leq M_i + \frac{\varepsilon}{4(b-a)}, \quad m_i^\varphi \geq m_i - \frac{\varepsilon}{4(b-a)}.
\]

Таким образом, колебание $\varphi(x)$ на $[x_i, x_{i+1}]$:
\[
\omega_i^\varphi = M_i^\varphi - m_i^\varphi \leq \left(M_i + \frac{\varepsilon}{4(b-a)}\right) - \left(m_i - \frac{\varepsilon}{4(b-a)}\right) = \omega_i + \frac{\varepsilon}{2(b-a)}.
\]

Теперь оценим сумму колебаний $\varphi(x)$ по всему разбиению:
\[
\sum_{i=0}^{n-1} \omega_i^\varphi \Delta x_i \leq \sum_{i=0}^{n-1} \left( \omega_i + \frac{\varepsilon}{2(b-a)} \right) \Delta x_i = \sum_{i=0}^{n-1} \omega_i \Delta x_i + \frac{\varepsilon}{2(b-a)} \sum_{i=0}^{n-1} \Delta x_i < \frac{\varepsilon}{2} + \frac{\varepsilon}{2(b-a)} \cdot (b - a) = \frac{\varepsilon}{2} + \frac{\varepsilon}{2} = \varepsilon.
\]

Следовательно, $\varphi(x)$ интегрируема по Риману.

\end{document}
